\section{Support-sensitive alternating conjugate generation}
\label{sec:SF-support-conjugate-generation}

\paragraph{Author freeze and origin.}
D-TD30--31, reserved question Q2eo, is a derivation by
\texttt{/root/sixfolds}. The commutator reduction and explicit constant
below are new derivations in this run. The generation step specializes
and fully reconstructs Hall's published Corollary 2.2; the small-support
case invokes the published Guralnick--Saxl Lemma 6.1. The spectral-rank
application was proposed by \texttt{/root} and reconstructed here.
The complete author proof, source-interface, checker and both output
readings and the root and independent whole-proof reviews pass.
No quotient is asserted to lift to a subgroup.

For $n\ge5$ and $1\ne g\in S_n$, let $k=|\operatorname{supp}(g)|$ and
let $\alpha_n(g)$ be the minimum number of $A_n$-conjugates of $g$
that generate $\langle A_n,g\rangle$. Thus this is generation by
conjugates, not a bound on word length in an entire conjugacy class.

\subsection*{D-TD30. A uniform bound depending on natural support}

\begin{theorem}[D-TD30]
For every $n\ge5$ and $1\ne g\in S_n$,
\[
                  \alpha_n(g)\le\frac{24n}{k}.
                  \tag{SCG.1}
\]
If $k\ge24$, the proof gives the stronger bound
$\alpha_n(g)\le16n/k$. Both assertions use $A_n$-conjugation,
including when $g$ is odd.
\end{theorem}

\begin{proof}
We first prove the needed generation lemma for a product of $s$
disjoint $3$-cycles, and then construct that product from $g$.

\paragraph{Step 1: a published generation lemma, with full reconstruction.}
Suppose $s\ge2$ and $n\ge6s-1$, and let $y\in A_n$ have cycle type
$3^s1^{n-3s}$. We claim
\[
        \alpha_n(y)\le 2+\left\lceil\frac{n-2}{2s}\right\rceil.
        \tag{SCG.2}
\]
This is the $p=3$ case of \cite[Corollary2.2, author copy p.4]{SF:Hall1995Finitary}.
We supply every step, including the transitive subgroup construction.

First, any two permutations of this cycle type are conjugate in $A_n$.
They are conjugate in $S_n$, and the centralizer of either contains an
odd permutation interchanging two of its three-cycles: in matching
cyclic coordinates that permutation is the product of three
transpositions. Composing an odd conjugator with it makes the
conjugator even. We may consequently choose any permutations of the
displayed cycle type in what follows.

Make a path with $2s$ vertices, colored alternately red and blue.
Associate a triple of letters to each vertex, so consecutive triples
intersect in one letter, every such intersection is a different letter,
and nonconsecutive triples are disjoint. The endpoint triples each
have two additional letters, and every other triple has one. There
are exactly $4s+1$ letters. Triples of either one color are disjoint.
The products of the $s$ three-cycles on the red and on the blue
triples are two conjugates of $y$. Their generated group is transitive
on these $4s+1$ letters, because the orbit equivalence relation is
generated by the triple incidences. They fix all other letters.

We place these letters inside a set $U_*$ of size $n-1$.
This is possible because $n-1\ge6s-2\ge4s+1$ for $s\ge2$.
If an already generated subgroup is transitive on $U\subseteq U_*$,
with $|U|\ge4s+1$, a further conjugate of $y$ can adjoin any prescribed
$q$ new letters, where $1\le q\le\min(2s,|U_*\setminus U|)$.
Write $q=2a+b$, with $b\in\{0,1\}$ and $a+b\le s$.
Use $a$ triples with two new and one old letter, $b$ triples with
one new and two old letters, and $s-a-b$ triples entirely in $U$.
The number of distinct old letters needed is $3s-q\le3s$;
there are enough, and all triples can be chosen disjoint.
Each newly added letter is joined to the existing orbit, so the
new subgroup is transitive on the enlarged set and still fixes
its complement.

The number of generators needed to reach $U_*$ is therefore
\[
 2+\left\lceil\frac{n-1-(4s+1)}{2s}\right\rceil
       =\left\lceil\frac{n-2}{2s}\right\rceil.
       \tag{SCG.3}
\]
Denote the one remaining letter by $z$. Add one conjugate of $y$
which moves $z$. The resulting group is transitive, and its stabilizer
of $z$ contains the preceding group transitive on $U_*$. It is thus
doubly transitive, and in particular primitive.

Add one more conjugate $y'$ of the first selected $y$ whose support
intersects that of $y$ in exactly one letter. The support union needs
$6s-1$ letters, so the hypothesis guarantees this choice. Only the
two three-cycles through the common letter fail to commute.
Writing these cycles as $(a\ b\ c)$ and $(a\ d\ e)$, direct
composition shows that their commutator is a single three-cycle.
Thus the enlarged primitive group contains a three-cycle.

For completeness, the required elementary primitive-group fact is
\cite[Proposition2.1 and proof, author copy p.4]{SF:Hall1995Finitary}.
Here is a proof. Two distinct three-cycles with overlapping supports
generate the alternating group on the union of their supports.
If the supports coincide, the assertion is simply that a three-cycle
generates $A_3$. For intersection size two, their transitive subgroup on four letters
has order divisible by both four and three, hence is $A_4$.
For intersection size one, writing them as $(a\ b\ c)$ and
$(a\ d\ e)$, conjugation gives $(b\ d\ e)$; these last two
cycles give $A_{\{a,b,d,e\}}$ by the preceding case. Conjugating
the first cycle by that $A_4$ gives every three-cycle containing
$c$ and two letters from $\{a,b,d,e\}$, and hence gives $A_5$.
This argument also proves inductively that three-cycles with a
connected support hypergraph generate the alternating group on
its vertices: merge a new support with the previously generated
alternating group, using one of these two overlap cases and then
conjugating by the old alternating group.

In a primitive group containing a three-cycle, take the hypergraph
of the supports of all its group conjugates. Its connected
components form an invariant block partition. Some component has
at least three vertices, so primitivity forces a single component.
The previous argument proves that the group contains $A_n$.
Our constructed group lies in $A_n$, and is therefore exactly $A_n$.
The two additions establish (SCG.2).

\paragraph{Step 2: even commutators supported in disjoint local blocks.}
Write the cycle decomposition of $g$. Let $a_2$ be its number of
2-cycles, let $k_{\ge3}$ be the sum of the lengths of its cycles
of length at least three, and put
\[
 m=\sum_{\ell\ge3\text{ cycles of }g}\left\lfloor\frac\ell3\right\rfloor,
 \qquad
 s_0=2\left\lfloor\frac m2\right\rfloor
               +2\left\lfloor\frac{a_2}3\right\rfloor,
 \qquad k=k_{\ge3}+2a_2.
 \tag{SCG.4}
\]
On a cycle $(a_0\ a_1\ \cdots\ a_{\ell-1})$, use the disjoint
consecutive triples
$(a_{3j},a_{3j+1},a_{3j+2})$, for
$0\le j<\lfloor\ell/3\rfloor$.
The adjacent transposition $t_j=(a_{3j}\ a_{3j+1})$ has conjugate
$g t_jg^{-1}=(a_{3j+1}\ a_{3j+2})$.
Their product is a three-cycle on that triple. Any chosen even
number of these transpositions, across all cycles, is an even
permutation; its commutator with $g$ is the product of the
corresponding disjoint three-cycles.

On three distinct 2-cycles $(u_i\ v_i)$, $i=1,2,3$, take the even
permutation $h_0=(u_1\ u_2\ u_3)$. Then
\[
        g h_0g^{-1}h_0^{-1}
          =(v_1\ v_2\ v_3)(u_1\ u_3\ u_2).
          \tag{SCG.5}
\]
It is a product of two disjoint three-cycles. Different groups of
three 2-cycles, and all preceding longer-cycle blocks, are disjoint
both before and after conjugation by $g$.

Pair the selected adjacent transpositions into units, and regard
each triple of 2-cycles as one further unit. Every unit uses an
even permutation and produces two disjoint three-cycles. Thus for
every even integer $s$ with $0\le s\le s_0$, one may select $s/2$
units and form an even permutation $h\in A_n$ such that
\[
                 [g,h]=g h g^{-1}h^{-1}
                   \quad\hbox{has cycle type }3^s1^{n-3s}.
                 \tag{SCG.6}
\]
This statement does not need any parity adjustment in a centralizer
of $g$, including when its $S_n$ conjugacy class splits in $A_n$.

\paragraph{Step 3: all endpoint inequalities for $k\ge24$.}
For every integer $\ell\ge3$,
$\lfloor\ell/3\rfloor\ge\ell/5$; writing $\ell=3q+r$ proves this
from $q\ge1$ and $0\le r\le2$. Also
$2\lfloor m/2\rfloor\ge m-1$ and
$2\lfloor a_2/3\rfloor\ge2a_2/3-4/3$.
Consequently
\[
 s_0\ge\frac{k_{\ge3}}5-1+\frac{2a_2}3-\frac43
       \ge\frac k5-\frac73\ge\frac{k}{10}
                       \qquad(k\ge24).
 \tag{SCG.7}
\]
Choose the even integer
\[
             s=\min\left\{s_0,\ 2\left\lfloor\frac{n+1}{12}\right\rfloor\right\}.
             \tag{SCG.8}
\]
To check the second entry is at least $n/10$, write
$n+1=12q+r$, with $0\le r\le11$.
Since $n\ge k\ge24$, $q\ge2$, and
$n\le12q+10\le20q$. Thus $2q\ge n/10$.
It follows that $s\ge k/10$, and $s$ is positive and even
(in fact $s\ge4$). By its upper cap, $6s\le n+1$.
All hypotheses of Step 1 hold.

Take the even $h$ of Step 2 and put $y=[g,h]$.
Each $A_n$-conjugate of $y$ lies in a subgroup generated by two
$A_n$-conjugates of $g$: for $c\in A_n$,
\[
 c y c^{-1}
    =(c g c^{-1})\bigl((c h)g(c h)^{-1}\bigr)^{-1}.
    \tag{SCG.9}
\]
Here both $c$ and $ch$ are even. Select the conjugates of $y$
that generate $A_n$ using Step 1 and collect the corresponding
conjugates of $g$. Their group contains $A_n$. If $g$ is even it
is exactly $A_n$, and if $g$ is odd it contains an odd generator
and is exactly $S_n$. Therefore
\[
 \begin{split}
 \alpha_n(g)&\le2\alpha_n(y)
       \le4+2\left\lceil\frac{n-2}{2s}\right\rceil\\
       &\le6+\frac{n-2}{s}
        \le6+\frac{10n}{k}\le\frac{16n}{k}.
 \end{split}
 \tag{SCG.10}
\]
The last step uses $n/k\ge1$.

\paragraph{Step 4: small support.}
For $2\le k\le23$, the published bound
\cite[Lemma6.1, pp.547--548]{SF:GuralnickSaxl2003Conjugates}
uses exactly our $A_n$-conjugation convention and gives
$\alpha_n(g)\le n-1$ for every $n\ge5$.
In particular $\alpha_n(g)<24n/k$.
Together with Step 3 this proves (SCG.1).
\end{proof}

\subsection*{D-TD31. Spectral rank on an actual perfect central cover}

\begin{theorem}[D-TD31]
Let $E$ be an actual finite perfect central extension of $A_n$,
$n\ge5$, and let $B$ be a nontrivial irreducible representation of
$E$ of dimension $r$ over an algebraically closed field of arbitrary
characteristic. For $e\in E$ whose image $y\in A_n$ moves $k>0$
letters, define the projective spectral rank by
\[
           \nu_B(e)=r-\max_{\lambda}\dim\ker(B(e)-\lambda I).
           \tag{SCG.11}
\]
Then
\[
                       \nu_B(e)\ge\frac{r k}{24n}.
                       \tag{SCG.12}
\]
No semisimplicity assumption is made on the individual matrix $B(e)$.
\end{theorem}

\begin{proof}
Take $c\le24n/k$ conjugates of $y$ under $A_n$ which generate
$A_n$, and lift their conjugating elements to $E$. The corresponding
conjugates $e_1,\ldots,e_c$ of $e$ generate a subgroup $H$ whose
image is $A_n$. If $Z$ denotes the central kernel, then $E=HZ$.
Perfectness gives
\[
                  E=[E,E]=[HZ,HZ]=[H,H]\le H,
                  \tag{SCG.13}
\]
so these actual conjugates generate $E$ itself.

Choose an eigenvalue $\lambda$ for which the eigenspace of $B(e)$
has maximum dimension. All the conjugates have an eigenspace of
that same dimension for $\lambda$. The intersection of those
$c$ eigenspaces has dimension at least $r-c\nu_B(e)$.
If this number were positive, a nonzero vector in the intersection
would span a line stable under the generated group $E$.
The action on that line would be a one-dimensional character of
a perfect group, hence trivial. A nontrivial irreducible
representation of a perfect group has no such line. Thus
$c\nu_B(e)\ge r$, and (SCG.12) follows.
The dimension-intersection argument works for genuine eigenspaces
in every characteristic and does not diagonalize $B(e)$.
\end{proof}

\paragraph{Precisely stated odd-extension variant.}
The same bound holds if an actual finite group $G=\langle E,x\rangle$
has a central-kernel surjection $\pi:G\longrightarrow S_n$,
$\pi(E)=A_n$, $E$ is perfect, and $\pi(x)$ is odd, provided the
representation of $G$ restricts irreducibly and nontrivially to $E$.
Indeed the $E$-conjugates of $x$ provided by D-TD30 generate a
subgroup $H$ with $G=H\ker\pi$. The normal subgroup
$\pi^{-1}(A_n)=E\ker\pi$ has derived subgroup $E$, so $E$
is normal in $G$. Also $G/E$ is cyclic, so
$G'=E$ by perfectness. Centrality gives $E=G'\le H$.
One of the selected conjugates is $a x a^{-1}$ with $a\in E\le H$,
whence $x\in H$ and $G=H$. The same common-eigenspace argument
produces a $G$-stable line whose restriction to $E$ contradicts
the stated irreducibility. This variant is not asserted for an
arbitrary semilinear normalizer or an abstract alternating quotient:
the actual central-kernel extension and the representation extension
are explicit additional hypotheses.

\paragraph{Source and application boundary.}
Hall's published chapter is \emph{Locally Finite Simple Groups of
Finitary Linear Transformations}, in \emph{Finite and Locally Finite
Groups}, Kluwer 1995, pp.147--188. The author-hosted copy has
Corollary2.2 and the complete Proposition2.1 proof on PDF page4;
its unresolved internal bibliography placeholders do not enter
the reconstruction above. Publisher and author metadata independently
identify this as a published chapter. The arbitrary-support constant
24 is derived here, not attributed as a stated Hall or
Hall--Liebeck--Seitz theorem. Guralnick--Saxl's published Lemma6.1
supplies only the coarse small-support bound. The stronger upstream
Hall--Liebeck--Seitz Theorem4 is not assumed: its full text has
not been recovered in this search.

D-TD31 applies to an already established actual perfect central
cover and an actual representation. It neither constructs that
cover from a quotient nor proves a minimum degree, tensor transfer,
Heisenberg trace bound, period rank, determinant comparison, or
new geometric source. Each such interface must be established
separately before applying these two group-theoretic results.
